\paragraph{纤维谱的有理完全性与交换对称 MOM 坐标：有限重建、统一逼近、零温集中与 Euler 缺陷能量}
在固定分辨率 \(m\) 的层面，本文的碰撞矩与 escort 机制不仅提供指数尺度的压力/谱读出，也给出一组严格的“有限证书化”接口：它们把纤维多重度的离散谱、交换对称可观测量的可学习空间以及 Euler 缺陷的交叉能量统一压缩为可计算且可审计的有限对象。其结论性接口可概括为：
\begin{enumerate}
  \item \textbf{纤维谱 resolvent 的有理结构与完全重建：有限个幂和矩决定全部极点与重数。}
  固定 \(m\) 后，幂和矩生成函数 \(H_m(u)=\sum_{q\ge 0}S_q(m)u^q\) 具有严格部分分式展开，其极点集合与留数分别编码纤维多重度的不同取值 \(\{a_j\}\) 与其重数 \(\{c_j\}\)（定理 \ref{thm:pom-fiber-spectrum-resolvent-rational}）。
  进一步，若 \(R\) 为该层的不同多重度取值个数，则至多 \(2R\) 个幂和数据 \(\{S_q(m)\}_{q=0}^{2R-1}\) 即可唯一重建 \(H_m\)，并且在一般情形下少一项将失去唯一性（定理 \ref{thm:pom-fiber-spectrum-finite-reconstruction-sharp}）。

  \item \textbf{交换对称可观测量的统一逼近：\texorpdfstring{$\{p_q\}$}{pq} 生成全部可学习对称统计量。}
  在概率单纯形 \(\Delta(X_m)\) 上，幂和特征 \(p_q(w)=\sum_x w(x)^q\)（\(2\le q\le |X_m|\)）生成的代数在一致范数下稠密于所有交换对称连续可观测量空间（定理 \ref{thm:pom-mom-symmetric-observables-stone-weierstrass}）。
  因而在“对类型标签不敏感”的学习/审计口径下，有限个碰撞矩特征已足以统一逼近任意对称统计量。

  \item \textbf{碰撞矩坐标图：一般位置下有限幂和给出对称分布空间的局部坐标。}
  在单纯形内点且权重两两不等的开稠密集合上，映射 \(w\mapsto(p_2(w),\dots,p_n(w))\) 的 Jacobian 由 Vandermonde 行列式控制，从而在交换对称意义下成为局部坐标图（定理 \ref{thm:pom-power-sum-vandermonde-jacobian}；推论 \ref{cor:pom-power-sum-local-chart}）。

  \item \textbf{零温极限定量律：escort 在 \texorpdfstring{$q\sim m$}{q~m} 标度下呈现 \texorpdfstring{$m^2$}{m^2} 级超指数集中。}
  escort 分布 \(w_{m,q}(x)\propto d_m(x)^q\) 对“低于最坏指数 \(\alpha_m-\varepsilon\)”的集合质量满足显式上界 \(w_{m,q}(A_m(\varepsilon))\le |X_m|e^{-q\varepsilon m}\)，从而当 \(q=\Theta(m)\) 时出现 \(\exp(-\Theta(m^2))\) 的超指数收缩（推论 \ref{cor:pom-escort-zero-temperature-superexp-concentration}）。

  \item \textbf{预言机反演的预算曲线：成功指数、容量指数与热相切割的闭式读出。}
  在 \(w_m\) 下，受限比特预算的最优反演成功率的指数主项由 \(I_w\) 通过折叠变分式严格给出，并可改写为双支极值（定理 \ref{thm:pom-oracle-success-variational-laplace}）。
  相应的容量指数满足 \(\sup_{\alpha}(f(\alpha)+\min(\alpha,a))\) 的变分表述，并在热相区间内由单层 \(\alpha=a\) 主导，从而多重分形计数谱可由容量曲线直接读出（定理 \ref{thm:pom-oracle-capacity-thermal-cut}；推论 \ref{cor:pom-oracle-capacity-slope-readout}）。
  同时微正则条件化与 escort 在层集上具有相对熵率为零的信息等价，临界预算窗满足高斯极限并给出二阶反演公式（定理 \ref{thm:pom-microcanonical-escort-kl-rate-zero}；定理 \ref{thm:pom-oracle-critical-window-gaussian}；推论 \ref{cor:pom-oracle-critical-window-bit-inversion}）。

  \item \textbf{Euler 缺陷的正交能量：交叉项 \texorpdfstring{$L^2$}{L2} 能量等于到可加子空间的最小距离。}
  在 \(\widehat G=\prod_i\widehat{G_i}\) 的均匀测度下，坐标平均算子诱导 Hoeffding 型正交分解并给出 Pythagoras 恒等式（定理 \ref{thm:pom-euler-defect-orthogonal-pythagoras}）。
  相应地，二阶及以上交叉分量的总能量严格等于到“常数 \(\,+\)”轴可加子空间 \(\mathcal S_{\mathrm{add}}\) 的最小 \(L^2\) 距离（定理 \ref{thm:pom-euler-defect-min-action-gauge}），从而提供一条规范化的 Euler 缺陷能量 gauge。
\end{enumerate}

